% Nouveau EGA V (cote 25): archive cote 25, site page 111 (PDF page 112).
% typescript. Draft compared: none.
% AI-transcribed by Claude Opus 5.5 (claude-opus-5-5, adaptive thinking, effort high), Message Batches API, batch ega25b (msgbatch_01A7n5YuXubUPNK1zsepi8RG), 2026-11-13; not yet checked by eye.
\scanpage{25/111}
\mspage{14}

la sedtion $\phi$ de $\underline{L}\otimes\underline{M}$ définie par \unsure{$\psi$} : $Q\longrightarrow\underline{M}$ soit transversalement
régulière.

Supposons maintenant que l'hypothèse $\underline{O}_S\xrightarrow{\sim}f_*(\underline{O}_X)$ soit encore vrai
après tout changement de base, ou ce qui revient au même grâce à III 7,
qu'on ait la condition
(xxx) \qquad\qquad $k(s)\xrightarrow{\sim}H^o(X_s,\underline{O}_{X_s})$
pour tout $s\in S$. Alors 2.12. s'applique également à tout $X_{S'}/S'$, pour un
changement de base quelconque $S'\rightarrow S$. On obtient ainsi le
\uline{Théorème} 2.13. Soit $f\colon X\longrightarrow S$ un morphisme \xout{plat} de présentation finie,
propre et plat satisfaisant à (xxx) ci-dessus, $\underline{L}$ un Module inversible
sur $X$, considérons le sous-foncteur $\mathrm{Div}^{L}_{X/S}$ de $\mathrm{Div}^{+}_{X/S}$ défini ci-dessus en
termes de $\underline{L}$. \xout{\unsure{Ce foncteur}} Il existe un Module de présentation finie $\underline{Q}$ sur
$S$ tel que le foncteur précédent soit représentable par un sous-préschéma
ouvert rétrocompact du fibré projectif $\underline{P}(\underline{Q})$ ; lorsque les fibres de $X/S$
sont géométriquement intègres, alors $\mathrm{Div}^{L}_{X/S}$ est représentable par ce
\inserted[typed]{fibré} \xout{\unsure{fibré}} projectif lui-même.

La dernière assertion résulte immédiatement de ce qu'on avait dit
avant. Pour le cas général, on avait déjà noté qu'on a un monomorphisme
\[
\mathrm{Div}^{L}_{X/S}\longrightarrow\underline{P}(\underline{Q}) \quad ,
\]
on est ramené à prouver que ce dernier est représentable par une immersion
ouverte quasi-compacte, ce qui nous ramène à prouver que si on se donne
une section de $P(\underline{Q})$ sur $S$ i.e. un Module quotient inversible $\underline{M}$ de $\underline{Q}$ ,
d'où une section $\phi$ de $\underline{L}\otimes\underline{M}$ non nulle sur chaque fibre, alors le \inserted{sous-}foncteur
$\underline{U}$ du foncteur final $S$ sur $(\mathrm{Sch})_{/S}$ ``consistant à rendre $\phi$ transversale-
ment régulière'' est représentable par un ouvert rétrocompact de $S$. Or
le fait qu'il est représentable par un ouvert \unsure{provieht|provient} du fait que $f$ est
propre et que la régularité transversale est une condition ouverte (cf par.
11 ...), la rétrocompacité se voit immédiatement par réduction au cas
noethérien.

\mspage{111}
